\author{Schmid Lothar}
\documentclass[12pt]{../../inc/mybib}

\setincpath{../../inc/}

\usepackage{bible_style}
\usepackage{textmakro}

\graphicspath{{../../assets/images/}}

% Absatzformatierung
\setlength{\parindent}{0pt}
\setlength{\parskip}{0.6em}

\begin{document}
\begin{bibelbox}{ESR}{Phil}{3:13-14}
    Brüder, ich selbst halte mich nicht dafür, [es] ergriffen zu haben; eines aber: Indem ich vergesse, was dahinten ist, und indem ich mich ausstrecke nach dem, was vorn ist,
    jage ich nach dem Ziel, hin zum Siegespreis, dem Ruf Gottes nach oben in Jesus, dem Gesalbten.
\end{bibelbox}
\subsection*{Aussage}
    Paulus weiss, dass er den Lauf noch nicht vollendet hat. Er will aber vorwärts gehen, ohne zurück zuschauen. Er will alles für den Herrn Jesus Christus einsetzten. 
\subsection*{Absicht}
    Paulus will seine Mitbrüder ermutigen. Er, der aktuell wegen Jesus Christus im Gefängnis sitzt, hätte das Recht zu verzweifeln. Aber nein er schaut nicht zurück. Er schaut nach vorne. Ausserdem will er sie daran Erinnern, dass nach dem Lauf auch ein Preis wartet. 
\subsection*{Anwendung}
    \textbf{Zielgruppe: Gemeindeleiter}\\
    Auch wenn nicht immer alles so verläuft wie man es geplant hat, soll man sich nicht entmutigen lassen. Es ist wichtig immer nach vorne zu schauen und sich am Wort Gottes zu richten. Man soll auch nicht das Ziel vergessen und der Lohn den uns Jesus Christus versprochen hat. Jesus ging ja uns voraus, um uns eine Wohnung zu schaffen. (\bibleverse{Joh}(14:2))
\subsection*{AIA}
    Schaue immer nach vorne und jage dem Ziel nach, denn so wirst du den Siegespreis erhalten.


\end{document}